\section*{S2. Data, force matching and prediction}

\subsection*{S2.1 Model and data}
The molecular dynamics model and simulation protocol are described in the main text. The additional interaction parameters are the FENE bond constants $K=30\epsilon/\sigma^2$ and $R_0=1.5\sigma$, the Lennard-Jones mixing rule $\sigma_{ij}=(\sigma_i+\sigma_j)/2$, and a $2\sigma$ cutoff for monomer pairs. Interactions involving ions are purely repulsive apart from electrostatics and the ion--monomer solvation potential \cite{tsamopoulos2024ion,shen2020ion},
\[
U_{\rm solv}(r)=-S\left[\left(\frac{\sigma}{r}\right)^4-\left(\frac{\sigma}{r_c}\right)^4\right],\qquad r\le r_c,
\]
with $S=4.33\epsilon$, $r_c=5\sigma$ and $\epsilon=k_BT$. Each concentration is equilibrated at zero pressure and sampled at its mean volume. The two dielectric constants use the same interaction parameters apart from electrostatics.
Four runs at each training concentration are held out, covering the different potential families, and three further runs are used for validation. 
Densities and internal force densities are recorded on $0.1\sigma$ bins. Profile uncertainties use eight time blocks. In the covariance-aware analysis, internal-force blocks are formed as $f_b^{\rm int}=f_b^{\rm tot}-f_b^{\rm ext}$ before computing their standard error.

\subsection*{S2.2 Force matching}
Let $\mathbf R$ denote all particle coordinates and
$\rmd\mathbf R$ the corresponding configurational measure.
The normalized canonical distribution is
\[
P(\mathbf R)
=
Z^{-1}\exp\left[
-\frac{1}{k_BT}
\left(
U_{\rm int}(\mathbf R)
+\sum_\alpha\sum_{j\in\alpha}V_\alpha(\vr_j)
\right)
\right],
\]
and the one-body density is
\[
n_\alpha(\vr)
=
\left\langle
\sum_{j\in\alpha}\delta(\vr-\vr_j)
\right\rangle
=
\sum_{j\in\alpha}\int\rmd\mathbf R\,
P(\mathbf R)\,\delta(\vr-\vr_j).
\]
Using
$\nabla_{\vr}\delta(\vr-\vr_j)
=-\nabla_j\delta(\vr-\vr_j)$
and integrating by parts over $\vr_j$, we obtain
\[
\begin{aligned}
\nabla n_\alpha(\vr)
&=
\sum_{j\in\alpha}\int\rmd\mathbf R\,
P(\mathbf R)\,\nabla_{\vr}\delta(\vr-\vr_j)\\
&=
-\sum_{j\in\alpha}\int\rmd\mathbf R\,
P(\mathbf R)\,\nabla_j\delta(\vr-\vr_j)\\
&=
\sum_{j\in\alpha}\int\rmd\mathbf R\,
\delta(\vr-\vr_j)\,\nabla_jP(\mathbf R).
\end{aligned}
\]
The boundary term vanishes for periodic boundaries
(with the periodic delta function), or for sufficiently
decaying distributions. For $j\in\alpha$,
\[
\nabla_jP(\mathbf R)
=
-\frac{P(\mathbf R)}{k_BT}
\left[
\nabla_jU_{\rm int}(\mathbf R)
+\nabla_jV_\alpha(\vr_j)
\right].
\]
Substitution therefore gives
\[
\begin{aligned}
k_BT\nabla n_\alpha(\vr)
={}&
-\left\langle
\sum_{j\in\alpha}\delta(\vr-\vr_j)
\nabla_jU_{\rm int}
\right\rangle\\
&-
\left\langle
\sum_{j\in\alpha}\delta(\vr-\vr_j)
\nabla_jV_\alpha(\vr_j)
\right\rangle.
\end{aligned}
\]
Since the delta function fixes $\vr_j=\vr$,
\[
\left\langle
\sum_{j\in\alpha}\delta(\vr-\vr_j)
\nabla_jV_\alpha(\vr_j)
\right\rangle
=
n_\alpha(\vr)\nabla V_\alpha(\vr).
\]
Defining the internal force density by
\[
\vf^{\rm int}_\alpha(\vr)
=
-\left\langle
\sum_{j\in\alpha}\delta(\vr-\vr_j)
\nabla_jU_{\rm int}
\right\rangle,
\]
we recover the local force balance
\cite{tschopp2022force,hansen2013theory},
\begin{equation}
k_BT\nabla n_\alpha(\vr)
=
\vf^{\rm int}_\alpha(\vr)
-n_\alpha(\vr)\nabla V_\alpha(\vr).
\label{eq:S:YBG}
\end{equation}
Taking the gradient of the Euler--Lagrange equation and using Eq.~\ref{eq:S:YBG} yields $\vf^{\rm int}_\alpha=-n_\alpha\nabla(e_\alpha\phi+\mu^{\rm ex}_\alpha)$, the force-matching relation of the main text. 

The six Gaussian envelope widths are $0.250$, $0.330$, $0.435$, $0.574$, $0.758$ and $1.000\sigma$. The pointwise network receives $h_{m\alpha}/n_{\rm ref}$, with $n_{\rm ref}=0.02\sigma^{-3}$, and its output contributes $n_{\rm ref}\int\Phi_\theta\,\rmd\vr$ to the free energy. Other architecture parameters are given in the main text. The kernel range and activation were selected using only
the validation profiles at $\varepsilon_r=2$.

Training uses full batches and the AdamW schedule specified in the main text, with learning rate reduced to $10^{-5}$. Validation is evaluated every 20 steps. Early stopping uses a patience of 1000 steps after at least 200 steps, and the best validation parameters are retained.
Table~\ref{tab:S:spectrum} compares training with and without the Fourier term, using the same architecture and data splits. The Fourier term improves the long-wavelength amplitude and number response and leaves the profile errors unchanged.

\begin{table}[H]
\centering
\caption{Loss ablation with and without the Fourier term weighted by $(k\sigma)^{-2}$. The amplitude ratio is the predicted over MD amplitude of the $k_1$ mode summed over the six long-wavelength training runs at $c=0.04$. At this concentration, MD gives $\Gamma(k_1)=4.38\pm0.15$ and $3.14\pm0.56$ at $\varepsilon_r=7.5$ and 2, respectively. In the last two columns, values before and after the slash
correspond to neutral ($\psi=0$) and mixed neutral/charged
($V_N\neq0$, $\psi\neq0$) external potentials, respectively.}
\label{tab:S:spectrum}
\begin{tabular}{c l c c c c c}
$\varepsilon_r$ & Fourier term & Ratio at $k_1$ & $\Gamma(k_1)$  & Profile (\%) & $k_1/2$ (\%) & $k_1/4$ (\%)\\
\midrule
7.5 & without & 0.95 & 4.56 & 1.45 & 18 / 8 & 13 / 8\\
7.5 & with & 0.96 & 4.47 & 1.50 & 16 / 6 & 11 / 6\\
2 & without & 0.87 & 3.85 & 3.09 & 11 / 10 & 15 / 17\\
2 & with & 0.92 & 3.70 & 3.11 & 7 / 7 & 12 / 14\\
\bottomrule
\end{tabular}
\end{table}

\subsection*{S2.3 Predictions and uncertainty}
At fixed particle numbers, the predicted planar profiles satisfy
\begin{equation}
n_\alpha(z)=N_\alpha\frac{e^{-u_\alpha(z)}}{A\int_0^L e^{-u_\alpha(z')}\,\rmd z'},\qquad
u_\alpha=V_\alpha+e_\alpha\phi[\vn]+\mu^{\rm ex}_\alpha[\vn],
\label{eq:S:canon}
\end{equation}
where $A$ is the cross-sectional area. The chemical potentials enforcing particle numbers drop out. We solve for $u_\alpha$ using Anderson acceleration with history eight and mixing parameter 0.1 \cite{anderson1965iterative,walker2011anderson}. The charge residual is preconditioned by $k^2/(k^2+\kappa_D^2)$ \cite{kerker1981efficient}. Convergence requires the maximum density change divided by $\nb$ to be below $10^{-6}$. Profile errors are $\|n_\alpha^{\rm pred}-n_\alpha^{\rm MD}\|_2/\|n_\alpha^{\rm MD}\|_2$, averaged over runs.

The bulk kernel $W_{\alpha\beta}(k)$ follows by automatic differentiation of $\mu^{\rm ex}_\alpha$ at uniform density along a one-bin perturbation of species $\beta$, followed by Fourier transformation. Equation~\ref{eq:S:OZ} gives $S(k)$. The neural functionals used for the main results have no negative eigenvalues of $S(k)^{-1}$ at the sampled concentrations and wavenumbers of a box eight times longer than the simulation box.

MD structure factors are obtained from zero-field runs.
For each frame, we compute the microscopic Fourier densities
\[
\tilde n_\alpha(\mathbf k,t)
=
\sum_{j\in\alpha}
e^{-i\mathbf k\cdot\vr_j(t)},
\qquad
\mathbf k=\frac{2\pi}{L}\mathbf m,
\quad \mathbf m\in\mathbb Z^3.
\]
For the homogeneous system at nonzero wavevectors,
$\langle\tilde n_\alpha(\mathbf k,t)\rangle=0$.
The number structure factor is therefore
\[
S_{NN}(k)
=
\frac{1}{V}
\left\langle
\left|
\tilde n_+(\mathbf k,t)+\tilde n_-(\mathbf k,t)
\right|^2
\right\rangle_{t,\,|\mathbf k|=k},
\]
where $V=L^3$ and the average is over time and wavevector
directions with the same magnitude.s
